\documentclass[a4paper, journal]{IEEEtran}
\IEEEoverridecommandlockouts

\usepackage{cite}
\usepackage{amsmath,amssymb,amsfonts}
\usepackage{algpseudocode}
\usepackage{graphicx}
\usepackage{textcomp}
\usepackage{listings}
\usepackage{xcolor}
\usepackage{booktabs}
\usepackage{tikz}
\usepackage{times}
\usepackage{newtxmath}
\usepackage{algorithm}
\renewcommand{\ttdefault}{cmtt}

\definecolor{codegreen}{rgb}{0,0.6,0}
\definecolor{codegray}{rgb}{0.2,0.2,0.2}
\definecolor{codepurple}{rgb}{0.58,0,0.82}
\definecolor{backcolour}{rgb}{0.95,0.95,0.92}

\lstset{
language=Python,
breaklines=true,
breakatwhitespace=true,
basicstyle=\ttfamily\footnotesize\color{codegray},
keywordstyle=\color{codepurple}\bfseries,
commentstyle=\color{codegreen}\itshape,
stringstyle=\color{red},
captionpos=b,
frame=shadowbox,
rulesepcolor=\color{gray},
showstringspaces=false,
tabsize=4,
numbers=left,
numberstyle=\tiny\color{codegray},
stepnumber=1,
numbersep=5pt,
backgroundcolor=\color{backcolour},
xleftmargin=15pt,
xrightmargin=15pt
}

\hyphenation{op-tical net-works semi-conduc-tor IEEE-Xplore}
\def\BibTeX{{\rm B\kern-.05em{\sc i\kern-.025em b}\kern-.08em
T\kern-.1667em\lower.7ex\hbox{E}\kern-.125emX}}
\usepackage{balance}
\begin{document}

\title{Family Relationship Inference System Using Python}

\author{
    \IEEEauthorblockN{Dinanath Padhya\IEEEauthorrefmark{1}, Krishna Acharya\IEEEauthorrefmark{2}} \\
    \IEEEauthorblockA{\IEEEauthorrefmark{1}THA078BEI010 \quad \IEEEauthorrefmark{2}THA078BEI020\\
        Department of Electronics, Communication and Information Engineering\\
        Thapathali Campus, IOE, Tribhuvan University, Nepal\\
    }}

\maketitle

\begin{abstract}
    This experiment presents a relationship inference engine using Python. The task includes adding family facts such as parent-child relations and allowing queries and listing of derived relationships like siblings, grandparents, uncles, and nieces. A rule-based approach using sets and tuples is employed to store known facts and compute additional relations via helper functions. The system supports various commands such as `add`, `list`, `query`, and `show`, with the aim of interpreting basic kinship from minimal input. This project not only tests the ability to manage simple data but also reinforces logical deduction based on established relationships. Through incremental development and structured logic, the program simulates a reasoning model that mirrors how genealogical relationships work in real life.
\end{abstract}

\section{Introduction}

Understanding and representing family relationships in a formalized way is a foundational problem in logic programming. The current lab is focused on implementing such a system using Python sets and tuples. Sets are used to ensure no duplicate relationships are stored, while tuples encode relationships such as (relation, person1, person2).

The lab implements a system that supports adding facts and inferring more complex relationships. For instance, if a parent relationship is defined, child and grandparent relations can be inferred using transitive rules. The implementation is command-driven, enabling operations such as querying specific relationships and listing all relatives of a given type.

Representing and processing relationships is one of the basic yet powerful tasks in computing, often used in the design of knowledge-based systems, simulations, and even in database architecture. Modeling such logic with basic programming constructs deepens understanding of the fundamentals of computer science, particularly in managing sets, tuple data structures, and algorithmic logic.

\section{Objectives}

\begin{itemize}
    \item To implement a basic fact storage system using Python sets.
    \item To infer complex relationships from basic facts.
    \item To enable queries and listings of supported relationships.
    \item To ensure the system works for both direct and indirect inferences.
    \item To practice modular design and logic flow within Python classes.
    \item To simulate rule-based relational deduction from minimal information.
    \item To reinforce tuple-based data organization and set operations.
\end{itemize}

\section{Supported Features}

The following features are implemented in the system:
\begin{itemize}
    \item Add parent or child relationships.
    \item Automatically infer reciprocal relationships.
    \item List relations such as parent, child, sibling, grandparent, grandchild, uncle/aunt, and nephew/niece.
    \item Query specific relationships between any two persons.
    \item View all stored facts.
    \item Dynamic computation of indirect relations based on existing known pairs.
    \item Prevention of redundant entries using set semantics.
\end{itemize}

\section{Methodology}

The core implementation revolves around a Python class named \texttt{FamilyRelations}. Facts are stored in a set as tuples in the form (relation, person1, person2). If a user adds a `child` relation, the reciprocal `parent` relation is automatically added. Each relationship is treated as a directed triple.

\subsection{Adding Facts}

The \texttt{add\_fact} method is used to add either a parent or child relation. If a child is added, the parent relation is inferred. Adding facts is fundamental to expanding the knowledge base of the program. By using consistent logic, even a small number of input relationships can yield a wide network of inferences.

\subsection{Retrieving Relationships}

Helper methods such as \texttt{get\_children}, \texttt{get\_parents}, \texttt{get\_siblings}, \texttt{get\_grandparents}, and others are used to retrieve relations from the stored facts. For example, siblings are identified by finding people with shared parents. Grandparents are found by tracing back two parent links.

\subsection{Inference}

Indirect relationships such as sibling or uncle are inferred using shared parents or sibling's children respectively. For example, siblings are identified by checking who shares the same parent. Grandchildren are derived from the children of a person's children. Uncles and aunts are derived by checking the siblings of a person's parents, and nieces or nephews from the children of a person's siblings.

\subsection{Query and List Commands}

\texttt{query(relation, person1, person2)}: Returns True if the specified relation holds between the two.

\texttt{list\_relation(relation, person)}: Lists all persons having the specified relation with the input person.

These commands help to verify and explore the relationship map interactively and ensure the model behaves as expected.

\subsection{Python Code Snippet}

\begin{lstlisting}[caption={Function to Get Children of a Person}]
def get_children(self, person):
return {child for (rel, p, child) in self.facts
if rel == 'parent' and p == person}
\end{lstlisting}

Other functions are structured similarly, iterating over the fact set and applying logic to filter the tuples based on the given relation and involved persons.

\section{Implementation Design}

The program is designed with reusability in mind. Each core relationship type has its own dedicated method. This modular design makes it easier to maintain and expand. It also improves readability and clarity. Edge cases, such as circular queries or unlisted persons, are handled gracefully with empty outputs.

For extensibility, the relation logic can be moved to a configuration-driven rule map, allowing users to define custom relationships without altering the source code.

\section{Testing and Evaluation}

Various relationships were added manually using the system and complex relations were correctly inferred. Some examples include:

\begin{itemize}
    \item Adding: parent Alice Bob
    \item Adding: parent Alice Charles
    \item Adding: parent Bob David
    \item Adding: parent Charles Emma
    \item Querying: sibling Bob Charles $\rightarrow$  Yes
    \item Querying: grandparent Alice Emma $\rightarrow$ Yes (via Charles)
    \item Listing: uncle Charles David $\rightarrow$ True
\end{itemize}

A sample of the program output when run:
\begin{verbatim}
All stored facts:
Parent: alice -> bob
Parent: alice -> charles
Parent: bob -> david
Parent: charles -> emma

List sibling bob: charles
List grandparent david: alice
Query parent alice bob: Yes
Query sibling bob charles: Yes
Query uncle charles david: Yes
\end{verbatim}

\section{Discussion}

The system accurately deduces complex relationships based on limited input. The use of sets eliminates duplicates and simplifies data management. Each command is stateless and derives results dynamically from existing facts.

Scenarios such as adopted relationships, step-siblings, or polygamous structures are not supported directly but can be added with minimal code changes by redefining the relationship logic.

The program encourages thinking in logical chains and understanding relational transitivity. Even though it is simple, the implications of its logic-based foundation are significant. Such a model can be extended toward decision support systems or as part of a broader genealogy software.

\section{Conclusion}

The experiment successfully implemented a basic relationship engine that can store and infer family relationships using only tuples and sets. The logic for indirect inference like siblings and grandparents was clearly verified. The approach ensures extensibility and simplicity by using fundamental Python data structures. This lab reinforces concepts in set theory, transitive relationships, and functional programming using Python.

The current structure is efficient for small to moderate datasets. Future versions may use indexing or graph data structures for faster querying in large-scale scenarios. Additional enhancements may include support for deleting facts, exporting relationships, and generating visual family trees from the facts.

\bibliographystyle{IEEEtran}
\bibliography{ref}

\end{document}
